\documentclass[11pt,reqno]{amsart}

\usepackage[a4paper,margin=31mm,headheight=14pt]{geometry}
\usepackage[T1]{fontenc}
\usepackage{newtxtext,newtxmath}
\usepackage{microtype}
\usepackage{amsmath,mathtools}
\usepackage{booktabs,longtable,array,tabularx}
\usepackage{enumitem}
\usepackage{needspace}
\usepackage{xurl}
\usepackage{csquotes}
\usepackage[backend=biber,style=numeric-comp,sorting=none,sortcites=true,maxbibnames=99,giveninits=true,doi=true,url=true,isbn=false]{biblatex}
\addbibresource{references.bib}
\AtBeginBibliography{\small}
\setlength{\bibitemsep}{0.14\baselineskip}
\usepackage[hidelinks]{hyperref}
\usepackage{bookmark}

\hypersetup{
  pdftitle={Obligatory Triple Systems: A Self-Contained Proof of Erdos Problem 593},
  pdfauthor={Samuil Petkov},
  pdfsubject={A self-contained ZFC proof of Erdos Problem 593}
}

\setlength{\parindent}{1.45em}
\setlength{\parskip}{0pt}
\setlist[itemize]{leftmargin=2.05em,itemsep=0.12em,topsep=0.32em}
\setlist[enumerate]{leftmargin=2.15em,itemsep=0.12em,topsep=0.32em}
\setcounter{tocdepth}{2}
\setcounter{secnumdepth}{0}
\renewcommand{\arraystretch}{1.12}
\setlength{\LTpre}{0.5em}
\setlength{\LTpost}{0.5em}
\emergencystretch=2.5em

\newcommand{\tightlist}{\setlength{\itemsep}{0pt}\setlength{\parskip}{0pt}}
\providecommand{\passthrough}[1]{#1}
\providecommand{\square}{\ensuremath{\Box}}

% Plain theorem typography: bold heading, italic statement, no colour or shadow.
\newenvironment{theorembox}[1]
  {\Needspace{8\baselineskip}\par\addvspace{0.78\baselineskip}\noindent\textbf{#1.}\enspace\itshape\ignorespaces}
  {\par\addvspace{0.34\baselineskip}}
\newenvironment{lemmabox}[1]
  {\Needspace{6\baselineskip}\par\addvspace{0.70\baselineskip}\noindent\textbf{#1.}\enspace\itshape\ignorespaces}
  {\par\addvspace{0.30\baselineskip}}
\newenvironment{propositionbox}[1]
  {\Needspace{6\baselineskip}\par\addvspace{0.70\baselineskip}\noindent\textbf{#1.}\enspace\itshape\ignorespaces}
  {\par\addvspace{0.30\baselineskip}}
\newenvironment{corollarybox}[1]
  {\Needspace{6\baselineskip}\par\addvspace{0.70\baselineskip}\noindent\textbf{#1.}\enspace\itshape\ignorespaces}
  {\par\addvspace{0.30\baselineskip}}
\newenvironment{claimbox}[1]
  {\Needspace{5\baselineskip}\par\addvspace{0.58\baselineskip}\noindent\textbf{#1.}\enspace\itshape\ignorespaces}
  {\par\addvspace{0.26\baselineskip}}
\newenvironment{definitionbox}[1]
  {\Needspace{6\baselineskip}\par\addvspace{0.70\baselineskip}\noindent\textbf{#1.}\enspace\normalfont\ignorespaces}
  {\par\addvspace{0.30\baselineskip}}
\newenvironment{remarkbox}[1]
  {\Needspace{5\baselineskip}\par\addvspace{0.58\baselineskip}\noindent\textit{#1.}\enspace\normalfont\ignorespaces}
  {\par\addvspace{0.26\baselineskip}}

\title[Obligatory Triple Systems]{Obligatory Triple Systems:\\A Self-Contained Proof of Erd\H{o}s Problem 593}
\author[Samuil Petkov]{Samuil Petkov}
\address{\'Ecole normale sup\'erieure -- PSL, Paris, France}
\email{samuil.petkov@ens.psl.eu}
\date{July 11, 2026}

\begin{document}

\begin{abstract}
We prove the exact characterisation of finite obligatory triple systems. After isolated vertices are removed, obligatoriness is equivalent both to membership in the class generated by private-vertex expansions of finite bipartite graphs under finite disjoint unions and one-point amalgamations, and to the intrinsic conjunction of linearity, one incident Levi bridge at every hyperedge-node, and evenness of every Berge cycle. The argument is valid in ZFC.
\end{abstract}

\maketitle

\begin{quote}\small
\textit{Bibliographic note.}
The author developed the results and arguments in this manuscript independently, without consulting Eric Li's manuscript, and learned of Li's contemporaneous preprint only afterward. Li also presents the same classification and, in a more general cardinal form, the same one-apex sequence lift, together with closely related Levi-bridge and quotient-forest arguments \autocite{li2026}. The present proofs do not use Li's results as premises; the citations record the contemporaneous work and the substantive overlap.
\end{quote}

\section{Conventions and definitions}\label{conventions-and-definitions}

A \emph{triple system} is a simple \(3\)-uniform hypergraph \(H=(V(H),E(H))\), so \(E(H)\subseteq[V(H)]^3\). A colouring \(c:V(H)\to\lambda\) is \emph{proper} when no hyperedge is monochromatic, and \(\chi(H)\) is the least cardinal \(\lambda\) admitting such a colouring. An embedding of a triple system \(F\) in \(H\) is an injective map \(f:V(F)\to V(H)\) such that \(f[e]\in E(H)\) for every \(e\in E(F)\); embeddings are not required to be induced. A finite triple system \(F\) is \emph{obligatory} if every triple system \(H\) with \(\chi(H)>\aleph_0\) contains an embedding of \(F\). A point is \emph{isolated} if it belongs to no hyperedge, and \(F^\circ\) denotes the result of deleting all isolated points.

All ordinary graphs in this paper are simple. A triple system is \emph{linear} when any two distinct hyperedges meet in at most one point. Its \emph{Levi graph} \(I(F)\) is the bipartite graph with classes \(V(F)\) and \(E(F)\), where a point-node \(p\) is adjacent to a hyperedge-node \(e\) exactly when \(p\in e\). A \emph{bridge} is an edge of a graph whose deletion increases the number of connected components. A \emph{Berge cycle of length} \(\ell\ge2\) consists of distinct points \(p_0,\ldots,p_{\ell-1}\) and distinct hyperedges \(e_0,\ldots,e_{\ell-1}\) such that \(p_i\in e_i\cap e_{i+1}\), with indices modulo \(\ell\). Equivalently, it is a simple cycle of length \(2\ell\) in \(I(F)\).

For a finite simple graph \(J\), its \emph{private-vertex expansion} \(J^+\) has vertex set
\[
V(J)\mathbin{\dot\cup}\{p_a:a\in E(J)\}
\]
and hyperedges \(\{u,v,p_{\{u,v\}}\}\) for \(\{u,v\}\in E(J)\), where the points \(p_a\) are new and pairwise distinct. A \emph{one-point amalgamation} of vertex-disjoint triple systems \(F_0,F_1\) is obtained by choosing \(x_i\in V(F_i)\), identifying \(x_0\) with \(x_1\), and making no other identifications or new hyperedges. Finite disjoint unions and finite edgeless systems include the empty cases.

\Needspace{10\baselineskip}
\section{Main theorem}\label{main-theorem}

\phantomsection
\label{theorem-a-exact-characterisation}
\begin{theorembox}{Theorem A: exact characterisation}

For every finite triple system \(F\), the following are equivalent:

\begin{enumerate}
\def\labelenumi{\arabic{enumi}.}
\tightlist
\item
  \(F\) is obligatory.
\item
  \(F\in\mathcal B\).
\item
  \(F^\circ\) is linear, every hyperedge-node of \(I(F^\circ)\) is incident with a bridge, and every Berge cycle of \(F^\circ\) has even length.
\end{enumerate}

Here \(F^\circ\) is obtained by deleting all isolated vertices, and \(\mathcal B\) is the smallest class containing every private-vertex expansion \(J^+\) of a finite bipartite graph, every finite edgeless triple system, and closed under finite disjoint unions and one-point amalgamations.

The proof is arranged as

\[
(2)\Longrightarrow(1),\qquad
(2)\Longleftrightarrow(3),\qquad
\neg(3)\Longrightarrow\neg(1).
\]

All arguments are in ZFC, and all embeddings are non-induced.

\end{theorembox}

\paragraph{Proof dependencies.}
The only non-elementary external results used as black boxes are the de Bruijn--Erd\H{o}s compactness theorem for finite graph colourings \autocite{debruijn1951} and the pair case of the Erd\H{o}s--Rado partition theorem \autocite{erdosrado1956}, namely
\[
(2^{\aleph_0})^+\longrightarrow(\aleph_1)^2_{\aleph_0}.
\]
Every other ingredient required for the classification is proved below.

\paragraph{Background.}
The problem was posed by Erd\H{o}s and is catalogued as Erd\H{o}s Problem 593 \autocite{erdos1995,bloom593}. Background on obligatory subsystems, large-chromatic set systems, and high-girth constructions may be found in \autocite{komjath2001,hajnalkomjath2008,ehr1973,erdoshajnal1966,egh1975,komjath2008,reiher2024,reihergirth2024,wang2026}.

\paragraph{Relation to Li's proof.}
Li's main theorem states the same classification, and Sections~3--5 of his preprint use the same one-apex sequence lift in a more general cardinal form, bridge selectors, and a quotient-forest reconstruction \autocite[Theorem~1.1 and Sections~3--5]{li2026}. The independently developed proof below instead proves the bipartite-expansion atom directly, deletes all Levi bridges in the intrinsic reconstruction, and packages the avoidance argument as an exact finite-linear-trace theorem. These are differences in proof and presentation, not a claim that the classification or all of its central structural ideas are new here.

\section{1. Preliminary lemmas}\label{preliminary-lemmas}

\phantomsection
\label{lemma-1.1-isolated-vertex-reduction}
\begin{lemmabox}{Lemma 1.1: isolated-vertex reduction}

Let \(F\) be a finite triple system and let \(F^\circ\) be obtained by deleting all isolated vertices. Then

\[
F\text{ is obligatory}\iff F^\circ\text{ is obligatory},
\]

and

\[
F\in\mathcal B\iff F^\circ\in\mathcal B.
\]

\end{lemmabox}
\textit{Proof.} If \(F^\circ\) is obligatory and \(H\) has uncountable chromatic number, then \(H\) has infinitely many vertices. After embedding the finite system \(F^\circ\), map the isolated vertices of \(F\) injectively to unused host vertices. Containment is non-induced, so additional host edges are irrelevant. The converse follows because \(F^\circ\) is a subhypergraph of \(F\).

For membership in \(\mathcal B\), the implication \(F^\circ\in\mathcal B\Rightarrow F\in\mathcal B\) follows by taking the disjoint union with a finite edgeless system.

For the converse, use structural induction on a construction of \(F\) in \(\mathcal B\). If \(F=J^+\), delete the isolated graph vertices of \(J\) to obtain \(J_0\); then \((J^+)^\circ=J_0^+\in\mathcal B\). If \(F\) is edgeless, its reduction is the empty edgeless system. Reduction commutes with disjoint union.

Suppose finally that \(F\) is a one-point amalgamation of \(F_0\) and \(F_1\), identifying \(x_i\in V(F_i)\). By induction \(F_0^\circ,F_1^\circ\in\mathcal B\). If both \(x_i\) are nonisolated in their factors, then \(F^\circ\) is their one-point amalgamation at those vertices. If exactly one \(x_i\) is nonisolated, the isolated side disappears at the identified point and \(F^\circ\) is the disjoint union of the two reductions. The same disjoint-union description applies when both \(x_i\) are isolated, because the identified point is deleted. Hence \(F^\circ\in\mathcal B\) in every case. \(\square\)

Thus isolated vertices may be suppressed throughout the substantive argument and restored at the end.

\phantomsection
\label{lemma-1.2-finite-deletion}
\begin{lemmabox}{Lemma 1.2: finite deletion}

If \(\chi(H)>\aleph_0\) and \(S\subseteq V(H)\) is finite, then

\[
\chi(H-S)>\aleph_0.
\]

\end{lemmabox}
\textit{Proof.} Otherwise a countable colouring of \(H-S\), together with finitely many fresh colours for \(S\), would colour \(H\) countably. \(\square\)

\phantomsection
\label{lemma-1.3-two-elementary-colouring-facts}
\begin{lemmabox}{Lemma 1.3: two elementary colouring facts}

The following facts will be used repeatedly.

\begin{enumerate}
\def\labelenumi{\arabic{enumi}.}
\tightlist
\item
  If \(E(H)=E_1\cup\cdots\cup E_r\) and every \((V(H),E_i)\) is countably chromatic, then \(H\) is countably chromatic.
\item
  If a graph \(G\) admits an orientation in which every vertex has outdegree at most \(d<\omega\), then \(\chi(G)\le 2d+1\).
\end{enumerate}

\end{lemmabox}
\textit{Proof.} For the first assertion, take the product of the \(r\) countable colourings.

For the second, every finite subgraph has at most \(d|V|\) edges and therefore a vertex of degree at most \(2d\). It is \((2d+1)\)-colourable by greedy deletion. The de Bruijn--Erdős compactness theorem \autocite{debruijn1951} then gives a \((2d+1)\)-colouring of all of \(G\). \(\square\)

\phantomsection
\label{lemma-1.4-closure-chain-lemma}
\begin{lemmabox}{Lemma 1.4: closure-chain lemma}

Let \(\kappa\) be an uncountable cardinal, let \(1\le r<\omega\), and let

\[
\Phi:[\kappa]^r\longrightarrow[\kappa]^{<\omega}
\]

have uniformly finite values. There is an increasing continuous chain

\[
\varnothing=M_0\subseteq M_1\subseteq\cdots
\subseteq M_i\subseteq\cdots\qquad(i<\operatorname{cf}\kappa)
\]

such that

\[
\bigcup_{i<\operatorname{cf}\kappa}M_i=\kappa,
\qquad |M_i|<\kappa,
\]

and every \(M_i\) is \(\Phi\)-closed:

\[
a\in[M_i]^r\quad\Longrightarrow\quad\Phi(a)\subseteq M_i.
\]

\end{lemmabox}
\textit{Proof.} Choose sets \(A_i\subseteq\kappa\), each of cardinality below \(\kappa\), whose union is \(\kappa\). At a successor stage close \(M_i\cup A_i\) under \(\Phi\), iterating the operation \(\omega\) times. Since \(\Phi\) is finitary and finite-valued, this does not increase an infinite cardinal below \(\kappa\). At a limit stage take the union. A union of an increasing chain of sets closed under a finitary operation remains closed; and because the stage is below \(\operatorname{cf}\kappa\), its cardinality remains below \(\kappa\). \(\square\)

Write

\[
I_i=M_{i+1}\setminus M_i.
\]

Then \((I_i)_{i<\operatorname{cf}\kappa}\) partitions \(\kappa\), and \(M_i=\bigcup_{j<i}I_j\).

\section{2. The graph lemma needed for the expansion atom}\label{the-graph-lemma-needed-for-the-expansion-atom}

\phantomsection
\label{lemma-2.1-uncountable-chromatic-number-forces-complete-bipartite-graphs}
\begin{lemmabox}{Lemma 2.1: uncountable chromatic number forces complete bipartite graphs}

For every positive integer \(n\), every graph of uncountable chromatic number contains \(K_{n,n}\).

\end{lemmabox}
\textit{Proof.} Suppose otherwise, and choose a \(K_{n,n}\)-free graph \(G\) of uncountable chromatic number whose vertex cardinality \(\kappa\) is minimal.

For \(A\in[\kappa]^n\), let

\[
N(A)=\{v\in\kappa:v\text{ is adjacent to every member of }A\}.
\]

Because \(G\) is \(K_{n,n}\)-free, \(|N(A)|<n\). Apply Lemma 1.4 to \(A\mapsto N(A)\), and let \(M_i,I_i\) be the resulting chain and layers.

By minimality of \(\kappa\), every induced graph \(G[I_i]\) is countably chromatic. Let \(G^\times\) consist of the edges joining distinct layers, and orient each such edge from the later endpoint toward the earlier endpoint. If \(v\in I_i\) had \(n\) earlier neighbours, those neighbours would form an \(n\)-set \(A\subseteq M_i\), and closure would give

\[
v\in N(A)\subseteq M_i,
\]

contrary to \(v\in I_i\). Thus the orientation of \(G^\times\) has outdegree at most \(n-1\). By Lemma 1.3, \(G^\times\) is finitely colourable.

Colour each \(G[I_i]\) with colours in \(\omega\), reusing the same palette on all layers, and combine this with a finite proper colouring of \(G^\times\). The resulting countable product colouring is proper on all of \(G\), a contradiction. \(\square\)

\section{3. Private-vertex expansions of bipartite graphs are obligatory}\label{private-vertex-expansions-of-bipartite-graphs-are-obligatory}

\phantomsection
\label{lemma-3.1-rainbow-bipartite-submatrices}
\begin{lemmabox}{Lemma 3.1: rainbow bipartite submatrices}

For all positive integers \(n,t\), there exists \(q=q(n,t)\) with the following property. Suppose the edges of \(K_{q,q}\) are coloured, with an arbitrary colour set, so that at every vertex each individual colour occurs on fewer than \(t\) incident edges. Then there are \(n\)-element subsets \(X'\) and \(Y'\) of the two vertex classes such that all \(n^2\) edges between \(X'\) and \(Y'\) have distinct colours.

\end{lemmabox}
\textit{Proof.} Choose independent uniformly random \(n\)-subsets \(X'\) and \(Y'\).

Let \(m_\gamma\) be the number of edges of colour \(\gamma\). The local bound gives

\[
m_\gamma\le(t-1)q.
\]

The number of pairs of same-coloured edges sharing a vertex is at most \((t-1)q^2\). The total number of pairs of same-coloured edges is at most

\[
\sum_\gamma\binom{m_\gamma}{2}
\le \frac12\bigl(\max_\gamma m_\gamma\bigr)\sum_\gamma m_\gamma
\le \frac12(t-1)q^3.
\]

For \(q\ge2n\), a fixed adjacent pair of edges is selected with probability at most \(2n^3/q^3\), while a fixed disjoint pair is selected with probability at most \(4n^4/q^4\). Hence the expected number of repeated colours in the selected \(K_{n,n}\) is at most

\[
\frac{2(t-1)n^3}{q}+\frac{2(t-1)n^4}{q}.
\]

This is below \(1\) for sufficiently large \(q\). Some choice of \(X',Y'\) therefore has no repeated colour. \(\square\)

\phantomsection
\label{proposition-3.2-the-complete-bipartite-expansion-atom}
\begin{propositionbox}{Proposition 3.2: the complete bipartite expansion atom}

For every positive integer \(n\), the triple system \(K_{n,n}^{+}\) is obligatory.

\end{propositionbox}
\textit{Proof.} Suppose not. Choose a \(K_{n,n}^{+}\)-free triple system

\[
H=(\kappa,E)
\]

with \(\chi(H)>\aleph_0\) and with \(\kappa\) minimal.

For a pair \(xy\), write

\[
d_H(xy)=|\{z:\{x,y,z\}\in E\}|.
\]

Set \(t=3n^2+1\) and define a graph \(R\) on \(\kappa\) by

\[
xy\in E(R)\quad\Longleftrightarrow\quad d_H(xy)\ge t.
\]

\subsubsection{\texorpdfstring{Claim 3.2.1: \(R\) is countably chromatic}{Claim 3.2.1: R is countably chromatic}}\label{claim-3.2.1-r-is-countably-chromatic}

Otherwise Lemma 2.1 gives a \(K_{n,n}\) in \(R\), with core vertex set \(C\), \(|C|=2n\). For each of its \(n^2\) graph edges \(xy\), choose a third vertex \(p_{xy}\) with

\[
\{x,y,p_{xy}\}\in E,
\]

avoiding \(C\) and all previously chosen private vertices. At every stage fewer than \(2n+n^2<t\) vertices are forbidden, while \(xy\) has at least \(t\) possible third vertices. Thus the \(p_{xy}\) can be chosen distinct and outside \(C\), producing a copy of \(K_{n,n}^{+}\), a contradiction. \(\square\)

Let

\[
E_1=\{e\in E:\text{some pair contained in }e\text{ is an edge of }R\},
\qquad E_2=E\setminus E_1.
\]

A proper colouring of \(R\) is a proper colouring of \((\kappa,E_1)\), since every \(E_1\)-edge contains an \(R\)-edge. Hence \(\chi(\kappa,E_1)\le\aleph_0\), and therefore

\[
\chi(\kappa,E_2)>\aleph_0. \tag{3.1}
\]

For a pair \(xy\), define

\[
\Phi(xy)=
\begin{cases}
\{z:\{x,y,z\}\in E\},&d_H(xy)<t,\\
\varnothing,&d_H(xy)\ge t.
\end{cases}
\]

Apply Lemma 1.4 to \(\Phi\), obtaining \(M_i,I_i\).

Every \(e\in E_2\) has at least two vertices in its highest layer. Indeed, suppose the highest layer met \(e\) only in \(z\in I_i\). Then the other pair \(xy=e\setminus\{z\}\) lies in \(M_i\). Since \(e\in E_2\), \(xy\notin R\), so \(d_H(xy)<t\). Closure gives \(z\in\Phi(xy)\subseteq M_i\), a contradiction.

Thus every \(E_2\)-edge is either contained in one layer, or has exactly two vertices in its highest layer \(I_i\) and its third vertex in \(M_i\). Let \(E_i^*\) be the \(E_2\)-edges contained in \(I_i\). Since \(|I_i|<\kappa\) and \(H\) is \(K_{n,n}^{+}\)-free, minimality of \(\kappa\) gives \(\chi(I_i,E_i^*)\le\aleph_0\). Consequently

\[
E^*=\bigcup_iE_i^*
\]

is countably chromatic. By (3.1), the remaining crossing-edge system

\[
E^{**}=E_2\setminus E^*
\]

has uncountable chromatic number.

For each \(i\), define a graph \(G_i\) on \(I_i\) by declaring \(xy\in E(G_i)\) when there is a \(\beta\in M_i\) such that \(\{x,y,\beta\}\in E^{**}\). If every \(G_i\) were countably chromatic, colour each \(I_i\) properly in \(G_i\), using the same countable palette on every layer. Every edge of \(E^{**}\) would then have differently coloured top-layer vertices, giving a countable colouring of \((\kappa,E^{**})\), a contradiction. Hence some \(G_i\) is uncountably chromatic.

Choose \(q=q(n,t)\) from Lemma 3.1. Lemma 2.1 gives a \(K_{q,q}\) in \(G_i\), with sides \(X,Y\). For every \(xy\in X\times Y\), choose

\[
\beta_{xy}\in M_i
\]

such that \(\{x,y,\beta_{xy}\}\in E^{**}\), and regard \(\beta_{xy}\) as the colour of the edge \(xy\).

This edge-colouring is locally \((t-1)\)-bounded. For example, if for fixed \(x\) and fixed \(\beta\) there were \(t\) different vertices \(y\) with \(\beta_{xy}=\beta\), then the pair \(\{x,\beta\}\) would have codegree at least \(t\). It would belong to \(R\), forcing the triples \(\{x,y,\beta\}\) into \(E_1\), contrary to their membership in \(E^{**}\subseteq E_2\). The same argument applies at vertices of \(Y\).

Lemma 3.1 gives \(n\)-sets \(X'\subseteq X\), \(Y'\subseteq Y\) for which the \(\beta_{xy}\), \(x\in X'\), \(y\in Y'\), are all distinct. They lie in \(M_i\), whereas \(X'\cup Y'\subseteq I_i\), so none is a core vertex. The triples

\[
\{x,y,\beta_{xy}\},
\qquad x\in X',\ y\in Y',
\]

form \(K_{n,n}^{+}\), the final contradiction. \(\square\)

\phantomsection
\label{corollary-3.3-every-finite-bipartite-expansion-is-obligatory}
\begin{corollarybox}{Corollary 3.3: every finite bipartite expansion is obligatory}

If \(J\) is any finite bipartite graph, then \(J^+\) is obligatory.

\end{corollarybox}
\textit{Proof.} Embed the two vertex classes of \(J\) in the two classes of some \(K_{n,n}\). Then \(J^+\) is a subhypergraph of \(K_{n,n}^+\). Obligatoriness is downward closed under non-induced containment: a copy of the larger finite system contains a copy of the smaller one. Apply Proposition 3.2. \(\square\)

\section{\texorpdfstring{4. Closure of obligatoriness under the operations defining \(\mathcal B\)}{4. Closure of obligatoriness under the operations defining \textbackslash mathcal B}}\label{closure-of-obligatoriness-under-the-operations-defining-mathcal-b}

\phantomsection
\label{lemma-4.1-disjoint-union-closure}
\begin{lemmabox}{Lemma 4.1: disjoint-union closure}

The class of finite obligatory triple systems is closed under finite disjoint unions.

\end{lemmabox}
\textit{Proof.} Let \(F_1,F_2\) be obligatory and let \(H\) be uncountably chromatic. First find a copy of \(F_1\). Delete its finite vertex set. By Lemma 1.2 the remaining triple system is still uncountably chromatic, so it contains \(F_2\). The two copies are disjoint. Iteration proves the finite case. \(\square\)

\phantomsection
\label{definition-4.2-rooted-abundance}
\begin{definitionbox}{Definition 4.2: rooted abundance}

Let \((F,r)\) be a finite rooted triple system with \(|V(F)|>1\). For a triple system \(H\) and \(m\ge1\), let \(R_m(F,r;H)\) be the set of vertices \(v\) for which there are \(m\) rooted copies of \(F\), all sending \(r\) to \(v\), whose off-root vertex sets are pairwise disjoint.

\end{definitionbox}
\phantomsection
\label{lemma-4.3-rooted-abundance-lemma}
\begin{lemmabox}{Lemma 4.3: rooted-abundance lemma}

If \((F,r)\) is finite and \(F\) is obligatory, and if \(\chi(H)>\aleph_0\), then

\[
\chi\bigl(H[R_m(F,r;H)]\bigr)>\aleph_0.
\]

\end{lemmabox}
\textit{Proof.} Put

\[
B=V(H)\setminus R_m(F,r;H).
\]

For each \(v\in B\), consider the family of off-root vertex sets of all copies of \(F\) rooted at \(v\). Its matching number is below \(m\). Choose a maximal pairwise disjoint subfamily and let \(S_v\) be the union of its members. Then

\[
|S_v|\le(m-1)(|V(F)|-1)=:D,
\]

and \(S_v\) meets the off-root part of every rooted copy at \(v\).

Form a graph \(D_B\) on \(B\) by joining \(v\) to each member of \(S_v\cap B\). If an edge is generated in both directions, assign it to one of its generators, and orient it away from that generator. Every vertex has outdegree at most \(D\), so Lemma 1.3 gives a finite proper colouring of \(D_B\).

Let \(A\) be one colour class. Then

\[
S_v\cap A=\varnothing\qquad(v\in A).
\]

Consequently \(H[A]\) contains no copy of \(F\). Indeed, if \(\varphi\) were such a copy and \(v=\varphi(r)\), its off-root part would have to meet \(S_v\) although both are subsets of \(A\). Since \(F\) is obligatory, every \(F\)-free triple system is countably chromatic. Thus each of the finitely many colour classes of \(D_B\) induces a countably chromatic subsystem. Tagging a vertex by its finite \(D_B\)-colour and its countable colour inside that class gives a countable proper colouring of \(H[B]\).

If \(H[R_m(F,r;H)]\) were countably chromatic as well, tagging the two sides would give a countable colouring of \(H\). This contradicts \(\chi(H)>\aleph_0\). \(\square\)

\phantomsection
\label{proposition-4.4-one-point-amalgamation-closure}
\begin{propositionbox}{Proposition 4.4: one-point amalgamation closure}

The class of finite obligatory triple systems is closed under one-point amalgamations.

\end{propositionbox}
\textit{Proof.} Let \(F_1,F_2\) be obligatory, with selected vertices \(r_1,r_2\), and let \(F\) be their one-point amalgamation. If one factor consists only of its selected vertex, the assertion is immediate. Otherwise set

\[
m=|V(F_2)|.
\]

By Lemma 4.3, the set

\[
R=R_m(F_1,r_1;H)
\]

induces an uncountably chromatic subsystem of every uncountably chromatic \(H\). Hence \(H[R]\) contains a copy of \(F_2\); write \(\varphi\) for its embedding and let \(v=\varphi(r_2)\).

At \(v\) there are \(m\) rooted copies of \(F_1\) with pairwise disjoint off-root sets. The set \(\varphi(V(F_2)\setminus\{r_2\})\) has \(m-1\) vertices, so it meets at most \(m-1\) of those off-root sets. One rooted \(F_1\)-copy therefore avoids the \(F_2\)-copy outside \(v\). Their union is the required one-point amalgamation. \(\square\)

Corollary 3.3, Lemma 4.1, Proposition 4.4, and the trivial edgeless case prove

\[
(2)\Longrightarrow(1). \tag{4.1}
\]

\section{5. Equivalence of the constructive and intrinsic descriptions}\label{equivalence-of-the-constructive-and-intrinsic-descriptions}

We prove \((2)\Longleftrightarrow(3)\). Isolated vertices may be ignored by Lemma 1.1.
A bridge-selector formulation and a related quotient-forest reconstruction appear in Li's proof \autocite[Sections~4--5]{li2026}; here we delete every Levi bridge and reconstruct directly from all bridge-deleted components.

\phantomsection
\label{proposition-5.1-the-intrinsic-conditions-are-preserved-by-the-generators-and-operations}
\begin{propositionbox}{Proposition 5.1: the intrinsic conditions are preserved by the generators and operations}

Every generator of \(\mathcal B\) satisfies the intrinsic conditions, and the conditions are preserved under finite disjoint unions and one-point amalgamations.

\end{propositionbox}
\textit{Proof.} For a bipartite graph \(J\), the expansion \(J^+\) is linear. Every hyperedge-node is joined to its private point by a bridge. A Berge cycle cannot use a private point, so the Berge cycles of \(J^+\) are precisely the ordinary graph cycles of \(J\), with the same lengths; these are even. Edgeless systems satisfy the conditions vacuously, and disjoint union plainly preserves them.

Consider a one-point amalgamation. Hyperedges belonging to different factors intersect only in the identified point, so linearity is preserved. Its Levi graph is the vertex-sum of the two Levi graphs at a point-node. Every simple cycle in a vertex-sum lies wholly in one factor: a cycle crossing into the other factor would have to pass through the cut vertex twice. Consequently bridges in the factors remain bridges, and every Berge cycle lies in one factor and retains its parity. \(\square\)

Thus

\[
(2)\Longrightarrow(3). \tag{5.1}
\]

\phantomsection
\label{proposition-5.2-bridge-block-decomposition}
\begin{propositionbox}{Proposition 5.2: bridge-block decomposition}

Suppose \(F\) is finite, has no isolated vertices, and satisfies the three intrinsic conditions. Then \(F\in\mathcal B\).

\end{propositionbox}
\textit{Proof.} If \(E(F)=\varnothing\), then the assumption that \(F\) has no isolated vertices gives \(V(F)=\varnothing\), and \(F\) is an edgeless generator. Hence assume \(E(F)\ne\varnothing\). Put \(L=I(F)\), let \(B(L)\) be its set of bridges, and let \(\mathscr C\) be the set of connected components of \(L-B(L)\).

Consider a hyperedge-node \(e\). Since \(e\) is incident with a bridge and has degree three, it has at most two nonbridge incidences. It cannot have exactly one: a nonbridge incidence lies on a cycle, and that cycle must leave \(e\) through a second nonbridge incidence. Therefore every hyperedge-node has either no nonbridge incidences or exactly two.

\subsubsection{The nontrivial components}\label{the-nontrivial-components}

Let \(C\in\mathscr C\) contain a nonbridge incidence. Every hyperedge-node \(e\in C\) has exactly two point-neighbours in \(C\). Contract each such \(e\) to an ordinary graph edge between those two point-neighbours. This produces a finite graph \(J_C\).

The graph \(J_C\) is simple: parallel edges would correspond to two hyperedges of \(F\) sharing two points, contrary to linearity. Every cycle in \(J_C\) gives a Berge cycle of the same length in \(F\). Hence \(J_C\) has no odd cycle and is bipartite.

Each hyperedge-node \(e\in C\) has a third point-neighbour \(p_e\) across a bridge. These third points satisfy:

\begin{enumerate}
\def\labelenumi{\arabic{enumi}.}
\tightlist
\item
  \(p_e\notin C\); otherwise a path in \(C\) from \(p_e\) to \(e\), together with the incidence \(ep_e\), would place that incidence on a cycle.
\item
  If \(e\ne f\) lie in \(C\), then \(p_e\ne p_f\); otherwise a path in \(C\) from \(e\) to \(f\), together with \(ep_e\) and \(p_ff\), would place those bridge incidences on a cycle.
\end{enumerate}

Thus the hyperedges belonging to \(C\), together with their third points \(p_e\), form exactly the private-vertex expansion \(J_C^+\).

\subsubsection{Components consisting of one hyperedge-node}\label{components-consisting-of-one-hyperedge-node}

If a hyperedge-node \(e\) has no nonbridge incidence, it is an isolated component after the bridges are deleted. Its single triple is the expansion of a one-edge bipartite graph. Thus every hyperedge of \(F\) belongs to a piece \(P_C\) isomorphic to \(J^+\) for a finite bipartite graph \(J\). The hyperedge sets of these pieces partition \(E(F)\). Their vertex sets cover \(V(F)\): every point is incident with a hyperedge, and it either lies in that hyperedge-node's bridge-deleted component or is the point endpoint of a bridge leaving it.

\subsubsection{Quotient forest}\label{quotient-forest}

Contract every component of \(L-B(L)\). The original bridges become edges of a quotient multigraph \(T\). Two distinct bridges cannot join the same pair of bridge-deleted components: paths inside those two components, together with the two bridges, would form a cycle containing both, contradicting that they are bridges. Thus \(T\) is a simple graph. It is a forest, because a cycle in the quotient would similarly lift to a cycle in \(L\) containing a purported bridge.

A vertex \(C\in V(T)=\mathscr C\) is called \emph{active} if the bridge-deleted component \(C\) contains a hyperedge-node. For an active \(C\), the piece \(P_C\) consists of the hyperedges represented in \(C\), all point-nodes of \(C\), and the point endpoints of bridge incidences leaving hyperedge-nodes in \(C\). An inactive component contains no hyperedge-node and hence no incidence of \(L-B(L)\), so it is a singleton point-node.

\phantomsection
\label{claim-5.2.1-running-intersection}
\begin{claimbox}{Claim 5.2.1: running intersection}

Every component of \(T\) contains an active vertex, because \(F\) has no isolated vertices. Choose an active root in each component, and order the active vertices by nondecreasing depth, with arbitrary order among equal-depth vertices. When an active component \(C\) is added, its piece \(P_C\) meets the union of the earlier pieces in at most one point.

\end{claimbox}
\textbf{Proof of the claim.} For a point \(p\in V(F)\), let \(X(p)\in\mathscr C\) be the component of \(L-B(L)\) containing the point-node \(p\). The active pieces containing \(p\) are precisely \(P_{X(p)}\), when \(X(p)\) is active, together with the pieces \(P_D\) at active neighbours \(D\) of \(X(p)\) whose corresponding original bridge has point endpoint \(p\). Thus the vertices of \(T\) indexing pieces that contain \(p\) are contained in the closed star centred at \(X(p)\).

Suppose that \(p\in P_C\) also lies in an earlier piece. If \(X(p)=C\), every other piece containing \(p\) is indexed by a neighbour of \(C\), and only the parent neighbour can be earlier. If \(X(p)\ne C\), then \(C\) is a neighbour of \(X(p)\). The vertex \(X(p)\) cannot be a child of \(C\), because then \(X(p)\) and all its other neighbours in the star would be deeper than \(C\). Hence \(X(p)\) is the parent of \(C\). In either case, \(p\) is the point endpoint of the unique parent edge of \(C\). The same reasoning applies to every point shared with an earlier piece, so all such points coincide with this parent point \(p_C\), and

\[
P_C\cap\bigcup_{D\text{ earlier}}P_D\subseteq\{p_C\}.
\]

For a root active component the intersection is empty. \(\square\)

Starting with the root piece in each component, add the active pieces in this order. If the new intersection is empty, use a disjoint union. If it is the singleton \(\{p_C\}\), use a one-point amalgamation at that point. Several earlier pieces may already contain the same inactive attachment point, but the new step still identifies only that one existing vertex. Because the pieces partition the hyperedges and cover the vertices of \(F\), the resulting triple system is exactly \(F\). Different components of \(T\) are combined by disjoint union. Hence \(F\in\mathcal B\). \(\square\)

This proves

\[
(3)\Longrightarrow(2). \tag{5.2}
\]

Combining (5.1) and (5.2),

\[
(2)\Longleftrightarrow(3). \tag{5.3}
\]

\section{6. The sequence lift used in the avoidance constructions}\label{the-sequence-lift-used-in-the-avoidance-constructions}

Li introduces the same one-apex sequence lift in the more general form \(\operatorname{Lift}(G,\kappa)\) and analyzes its bridge traces \autocite[Sections~3--4]{li2026}. Our construction is its \(\kappa=\omega_1\) specialization, followed here by a different finite-linear-trace decomposition; we give the complete argument.

Let \(G=(X,A)\) be an arbitrary graph, where \(A=E(G)\) is also regarded as an alphabet. Let

\[
T=A^{<\omega_1}
\]

be the tree of all sequences of graph edges of countable ordinal length. For \(s\in T\) and \(x\in X\), write \(x_s=(s,x)\).

\phantomsection
\label{definition-6.1-the-one-apex-lift}
\begin{definitionbox}{Definition 6.1: the one-apex lift}

Define the triple system \(\mathcal L(G)\) by

\[
V(\mathcal L(G))=T\times X.
\]

For \(s\in T\), an edge \(a=\{x,y\}\in A\), a sequence \(t\in T\) extending \(s^\frown a\), and \(z\in X\), put

\[
\{x_s,y_s,z_t\}\in E(\mathcal L(G)). \tag{6.1}
\]

The two vertices \(x_s,y_s\) form the base pair and \(z_t\) is the apex. Since \(t\) properly extends \(s\), these are three distinct vertices. In the resulting triple, the sequence node occurring twice uniquely identifies \(s\); its two graph labels identify the edge \(a=\{x,y\}\), and the remaining vertex identifies \(t,z\). Thus each hyperedge arises from unique data, up to interchanging \(x\) and \(y\), and the system is simple and \(3\)-uniform.

\end{definitionbox}
\phantomsection
\label{lemma-6.2-chromatic-lower-bound}
\begin{lemmabox}{Lemma 6.2: chromatic lower bound}

If \(\chi(G)>\aleph_0\), then

\[
\chi(\mathcal L(G))>\aleph_0.
\]

\end{lemmabox}
\textit{Proof.} Suppose \(c:T\times X\to\omega\) is proper. For every \(s\in T\), the restriction \(x\mapsto c(x_s)\) is a countable colouring of \(G\). Hence there is an edge

\[
a_s=\{x_s^0,x_s^1\}\in E(G)
\]

whose two vertices at node \(s\) have a common colour \(k_s\).

Construct \(s_\alpha\in A^\alpha\) for \(\alpha<\omega_1\) by

\[
s_0=\varnothing,
\qquad s_{\alpha+1}=s_\alpha^\frown a_{s_\alpha},
\]

and at limit ordinals take unions.

If \(\alpha<\beta\), then \(s_\beta\) extends \(s_\alpha^\frown a_{s_\alpha}\). If \(k_{s_\alpha}=k_{s_\beta}\), take either endpoint \(z\) of the monochromatic edge at \(s_\beta\). By (6.1),

\[
\{(s_\alpha,x_{s_\alpha}^0),
  (s_\alpha,x_{s_\alpha}^1),
  (s_\beta,z)\}
\]

is a monochromatic hyperedge. Therefore the colours \(k_{s_\alpha}\), \(\alpha<\omega_1\), are pairwise distinct, impossible in \(\omega\). \(\square\)

\phantomsection
\label{theorem-6.3-exact-finite-linear-trace-theorem}
\begin{theorembox}{Theorem 6.3: exact finite linear trace theorem}

Every finite linear subhypergraph \(K\) of \(\mathcal L(G)\) is obtainable, by finite disjoint unions and one-point amalgamations, from systems \(J^+\) where \(J\) is a finite subgraph of \(G\).

\end{theorembox}
\textit{Proof.} First set aside every isolated vertex of \(K\). Each is a one-vertex edgeless factor, hence isomorphic to \(J^+\) for a one-vertex edgeless subgraph \(J\) of \(G\), and may be restored by disjoint union at the end. We may therefore assume that every vertex of \(K\) lies in a hyperedge.

Every hyperedge of \(\mathcal L(G)\) has a unique base node \(s\): exactly two of its vertices have sequence coordinate \(s\), while the apex has a strictly longer sequence coordinate. Group the hyperedges of \(K\) by their base nodes, and let \(K_s\) be the subhypergraph whose hyperedges are those based at \(s\) and whose vertices are their union.

For a fixed graph edge \(a=\{x,y\}\in E(G)\), at most one hyperedge of \(K_s\) can have base pair \(\{x_s,y_s\}\), since two such hyperedges would share two vertices and violate linearity. Hyperedges of \(K_s\) associated with distinct graph edges have distinct apices: if \(z_t\) were the apex for graph edges \(a\) and \(a'\), the coordinate of \(t\) immediately after \(s\) would have to be both \(a\) and \(a'\). It follows that \(K_s\) is precisely \(J_s^+\), where \(J_s\) is the finite subgraph of \(G\) consisting of the graph edges used at \(s\).

Now consider intersections between different \(K_s\). If a point \(z_t\) lies in both \(K_s\) and \(K_u\), then both \(s\) and \(u\) are prefixes of \(t\), so they are comparable. Suppose \(s\subsetneq u\). A common point cannot be a core point of \(K_s\); it is an apex of \(K_s\). The graph edge used by its hyperedge in \(K_s\) is the coordinate of \(u\) immediately following \(s\). Hence all points in \(K_s\cap K_u\) would be apices of the same base edge of \(K_s\). There is at most one such apex. Therefore

\[
|V(K_s)\cap V(K_u)|\le1. \tag{6.2}
\]

Construct a bipartite incidence graph \(Q\) whose one class is the set of active base nodes \(s\), and whose other class consists of points lying in at least two of the \(K_s\); join \(s\) to \(p\) when \(p\in K_s\).

We claim that \(Q\) is a forest. Suppose it has a cycle, and choose on that cycle a base node \(s\) of minimum length. Its two neighbouring base nodes on the cycle are proper descendants of \(s\). The two intervening shared points are distinct, so they are distinct apices in \(K_s\), arising from two distinct graph edges \(a\ne a'\). Thus the two neighbouring base nodes lie in different immediate branches below \(s\), one beginning with \(a\) and the other with \(a'\).

On the remainder of the cycle, consecutive base nodes share a point and are therefore comparable. Every base node there has length at least that of \(s\). Comparable descendants of \(s\) lie in the same immediate branch below \(s\), so the remainder of the cycle cannot move from the \(a\)-branch to the \(a'\)-branch without passing through a base node shorter than \(s\). This contradicts the choice of \(s\). Hence \(Q\) is acyclic.

Root each component of \(Q\) at a base node, and order its base nodes by nondecreasing even distance from the root. For a nonroot base node \(s\), let \(p_s\) be the point-node adjacent to \(s\) on the path to the root. Since \(Q\) is a tree, any earlier factor meeting \(K_s\) does so at \(p_s\); earlier siblings can also meet \(K_s\), but only at this same point. Equation (6.2) then gives
\[
V(K_s)\cap\bigcup_{u\text{ earlier}}V(K_u)\subseteq\{p_s\}.
\]
Starting with the root factor, add the remaining factors in this order by disjoint unions or one-point amalgamations as appropriate. Different components of \(Q\), and the isolated factors removed at the start, are combined by disjoint union. \(\square\)

\phantomsection
\label{corollary-6.4-restrictions-on-finite-linear-traces}
\begin{corollarybox}{Corollary 6.4: restrictions on finite linear traces}

Every finite linear \(K\subseteq\mathcal L(G)\) satisfies:

\begin{enumerate}
\def\labelenumi{\arabic{enumi}.}
\tightlist
\item
  every hyperedge-node of \(I(K)\) is incident with a bridge;
\item
  every Berge cycle of \(K\) is contained in a single factor \(J_s^+\);
\item
  a Berge cycle of length \(\ell\) in \(K\) yields an ordinary cycle of length \(\ell\) in \(G\).
\end{enumerate}

\end{corollarybox}
\textit{Proof.} In every \(J_s^+\) the incidence with the private vertex is a bridge. Bridgehood is preserved by one-point amalgamation. A simple Levi cycle cannot cross a one-point cut, so every Berge cycle lies in one factor. Berge cycles in \(J_s^+\) are exactly ordinary graph cycles in \(J_s\). \(\square\)

\section{7. Uncountably chromatic graphs of large odd girth}\label{uncountably-chromatic-graphs-of-large-odd-girth}

For a positive integer \(r\) and an infinite cardinal \(\kappa\), define the directed shift graph \(D_r(\kappa)\) as follows:

\begin{itemize}
\tightlist
\item
  its vertices are increasing \(r\)-tuples \[(\alpha_0,\ldots,\alpha_{r-1}),\qquad
    \alpha_0<\cdots<\alpha_{r-1}<\kappa;\]
\item
  it has an arc \[(\alpha_0,\ldots,\alpha_{r-1})
    \longrightarrow
    (\alpha_1,\ldots,\alpha_r)\] whenever \(\alpha_0<\cdots<\alpha_r\).
\end{itemize}

Let \(S_r(\kappa)\) be its underlying graph. Write

\[
\beth_0(\aleph_0)=\aleph_0,
\qquad
\beth_{j+1}(\aleph_0)=2^{\beth_j(\aleph_0)}.
\]

\phantomsection
\label{lemma-7.1-chromatic-number-of-shift-graphs}
\begin{lemmabox}{Lemma 7.1: chromatic number of shift graphs}

If

\[
\kappa>\beth_{r-1}(\aleph_0),
\]

then

\[
\chi(S_r(\kappa))>\aleph_0.
\]

\end{lemmabox}
\textit{Proof.} The case \(r=1\) is immediate because \(S_1(\kappa)=K_\kappa\). Assume \(r\ge2\). Suppose \(S_r(\kappa)\) has a proper colouring \(c\) with \(\lambda\) colours. For every increasing \((r-1)\)-tuple \(u\), define

\[
C(u)=\{c(u^\frown\beta):\beta>\max u\}\subseteq\lambda.
\]

If

\[
u=(\alpha_0,\ldots,\alpha_{r-2}),
\qquad
v=(\alpha_1,\ldots,\alpha_{r-1})
\]

are adjacent in \(S_{r-1}(\kappa)\), then

\[
c(\alpha_0,\ldots,\alpha_{r-1})\in C(u).
\]

It cannot belong to \(C(v)\), since that would give some \(\beta>\alpha_{r-1}\) for which the adjacent \(r\)-tuples

\[
(\alpha_0,\ldots,\alpha_{r-1}),
\qquad
(\alpha_1,\ldots,\alpha_{r-1},\beta)
\]

have the same colour. Hence \(C(u)\ne C(v)\). Thus a \(\lambda\)-colouring of \(S_r(\kappa)\) yields a \(2^\lambda\)-colouring of \(S_{r-1}(\kappa)\).

Starting with a countable colouring of \(S_r(\kappa)\) and iterating \(r-1\) times would give a colouring of

\[
S_1(\kappa)=K_\kappa
\]

with at most \(\beth_{r-1}(\aleph_0)\) colours, impossible when \(\kappa\) is larger. \(\square\)

\phantomsection
\label{lemma-7.2-odd-girth-of-shift-graphs}
\begin{lemmabox}{Lemma 7.2: odd girth of shift graphs}

The odd girth of \(S_r(\kappa)\) is at least

\[
2r+1.
\]

\end{lemmabox}
\textit{Proof.} For a digraph \(D\), let \(\delta D\) be its arc digraph: the vertices of \(\delta D\) are the arcs of \(D\), and

\[
(u,v)\longrightarrow(v,w)
\]

is an arc of \(\delta D\). We have

\[
D_{r+1}(\kappa)\cong\delta D_r(\kappa).
\]

We use the following fact: if \(D\) is acyclic and the underlying graph of \(D\) has odd girth \(g\), then the underlying graph of \(\delta D\) has odd girth at least \(g+2\). Indeed, a directed cycle in \(\delta D\) would concatenate to a directed closed walk in \(D\), which would contain a directed cycle; hence \(\delta D\) is also acyclic.

Take an odd cycle \(a_0,\ldots,a_{m-1}\) in the underlying graph of \(\delta D\). Record, around the cycle, whether each adjacency is oriented forward or backward in \(\delta D\). Since \(\delta D\) is acyclic, the cyclic sign sequence is not constant. It therefore has at least two sign changes, and the number of sign changes is even.

For each adjacency \(a_{i-1},a_i\), take their common endpoint in \(D\). Consecutive such endpoints are joined by the arc \(a_i\) when the two adjacent signs agree, and are equal when the signs change. Deleting the stationary steps gives an odd closed walk in the underlying graph of \(D\), of length at most \(m-2\). Every odd closed walk contains an odd cycle, so \(m\ge g+2\).

Now \(D_1(\kappa)\) is the transitive tournament; it is acyclic and its underlying graph \(K_\kappa\) has odd girth \(3\). Iterating the preceding observation gives odd girth at least

\[
3+2(r-1)=2r+1.
\]

\(\square\)

\phantomsection
\label{corollary-7.3-zfc-graphs-with-uncountable-chromatic-number-and-prescribed-odd-girth}
\begin{corollarybox}{Corollary 7.3: ZFC graphs with uncountable chromatic number and prescribed odd girth}

For every finite \(m\) there is a graph \(G\) with \(\chi(G)>\aleph_0\) and no odd cycle of length at most \(m\).

\end{corollarybox}
\textit{Proof.} Choose \(r\) with \(2r+1>m\), and let

\[
G=S_r\bigl(\beth_{r-1}(\aleph_0)^+\bigr).
\]

Apply Lemmas 7.1 and 7.2. \(\square\)

\section{8. The avoidance hosts}\label{the-avoidance-hosts}

We prove

\[
\neg(3)\Longrightarrow\neg(1).
\]

\phantomsection
\label{proposition-8.1-avoidance-of-every-nonlinear-finite-triple-system}
\begin{propositionbox}{Proposition 8.1: avoidance of every nonlinear finite triple system}

There is an uncountably chromatic linear triple system. Consequently, every finite nonlinear triple system is non-obligatory.

\end{propositionbox}
\textit{Proof.} We use the standard Erdős--Rado relation \autocite{erdosrado1956}

\[
(2^{\aleph_0})^+\longrightarrow(\aleph_1)^2_{\aleph_0}. \tag{8.1}
\]

Let \(\kappa=(2^{\aleph_0})^+\). Define a triple system \(T_\kappa\) whose vertices are the pairs in \([\kappa]^2\), and whose hyperedges are

\[
\bigl\{
\{\alpha,\beta\},
\{\alpha,\gamma\},
\{\beta,\gamma\}
\bigr\},
\qquad \alpha<\beta<\gamma<\kappa.
\]

This triple system is linear: two distinct triangles of a complete graph cannot share two graph edges, because two graph edges of a triangle determine the third edge and the three underlying vertices.

It is uncountably chromatic. A countable vertex-colouring of \(T_\kappa\) is a countable edge-colouring of \(K_\kappa\). By (8.1) there is a monochromatic subset of cardinality \(\aleph_1\), and in particular a monochromatic graph triangle. This is a monochromatic hyperedge of \(T_\kappa\).

If a finite triple system \(F\) is nonlinear, it has two distinct edges sharing at least two vertices. No injective embedding of those edges can exist in the linear host \(T_\kappa\). Thus \(T_\kappa\) is \(F\)-free. \(\square\)

\phantomsection
\label{proposition-8.2-avoidance-when-the-bridge-condition-fails}
\begin{propositionbox}{Proposition 8.2: avoidance when the bridge condition fails}

Let \(F\) be finite and linear, and suppose some hyperedge-node of \(I(F^\circ)\) is incident with no bridge. Then \(F\) is non-obligatory.

\end{propositionbox}
\textit{Proof.} By Lemma 1.1 it is enough to avoid \(F^\circ\), so assume that \(F\) has no isolated vertices. Take \(G=K_{\omega_1}\). Then \(\chi(G)>\aleph_0\), and Lemma 6.2 gives

\[
\chi(\mathcal L(G))>\aleph_0.
\]

If \(F\) appeared in \(\mathcal L(G)\), retain exactly the host hyperedges corresponding to the edges of \(F\). Because the embedding is injective and \(F\) is linear, these selected host hyperedges form a finite linear subhypergraph isomorphic to \(F\). Corollary 6.4 says that every hyperedge-node of such a trace is incident with a bridge, contradicting the chosen hyperedge of \(F\). Therefore \(\mathcal L(K_{\omega_1})\) is an uncountably chromatic \(F\)-free host. \(\square\)

\phantomsection
\label{proposition-8.3-avoidance-of-an-odd-berge-cycle}
\begin{propositionbox}{Proposition 8.3: avoidance of an odd Berge cycle}

Let \(F\) be finite and linear, and suppose \(F^\circ\) has an odd Berge cycle. Then \(F\) is non-obligatory.

\end{propositionbox}
\textit{Proof.} Again suppress isolated vertices. Put \(m=|E(F)|\). Choose \(r\) with \(2r+1>m\), and set

\[
G=S_r\bigl(\beth_{r-1}(\aleph_0)^+\bigr).
\]

By Lemmas 7.1 and 7.2, \(\chi(G)>\aleph_0\), and \(G\) has no odd cycle of length at most \(m\). Lemma 6.2 gives

\[
\chi(\mathcal L(G))>\aleph_0.
\]

If \(F\) appeared in \(\mathcal L(G)\), the selected image hyperedges would form a finite linear trace. By Corollary 6.4, every Berge cycle of that trace lies in one factor \(J_s^+\) and induces an ordinary graph cycle of the same length in \(J_s\subseteq G\).

The odd Berge cycle of \(F\) has length at most \(m\), so \(G\) would contain an odd cycle of length at most \(m\), a contradiction. Thus \(\mathcal L(G)\) is an uncountably chromatic \(F\)-free host. \(\square\)

The three propositions cover every failure of the intrinsic conditions, and prove

\[
\neg(3)\Longrightarrow\neg(1). \tag{8.2}
\]

\section{9. Conclusion}\label{conclusion}

We have established:

\begin{itemize}
\tightlist
\item
  \((2)\Longrightarrow(1)\) from the direct bipartite-expansion atom and the two closure arguments;
\item
  \((2)\Longleftrightarrow(3)\) from the Levi bridge-block decomposition and running-intersection reconstruction;
\item
  \(\neg(3)\Longrightarrow\neg(1)\) from the explicit nonlinear host and the sequence-lift avoidance constructions.
\end{itemize}

Therefore

\begingroup\small
\[
\begin{gathered}
F\text{ is obligatory}
\iff
F\in\mathcal B,\\[1mm]
F\in\mathcal B
\iff
\begin{array}{l}
F^\circ\text{ is linear},\\[1mm]
\text{every hyperedge-node of }I(F^\circ)
\text{ is incident with a bridge},\\[1mm]
\text{every Berge cycle of }F^\circ\text{ has even length}.
\end{array}
\end{gathered}
\]
\endgroup

All embeddings are non-induced: additional host hyperedges on the image vertices are irrelevant throughout. The proof uses only ZFC, the displayed successor cardinals, graph-colouring compactness, and the standard Erdős--Rado relation (8.1). It assumes neither CH nor GCH, and no forcing axiom or large-cardinal hypothesis.


\clearpage
\printbibliography[heading=bibintoc,title={References}]

\section*{Ethics and AI-assistance statement}
This manuscript was completed with assistance from OpenAI's ChatGPT-5.6 Pro for proof development, adversarial checking, drafting, editing, bibliographic organisation, and typesetting. Samuil Petkov is the sole author and accepts responsibility for the verification of all mathematical arguments, citations, and conclusions. The AI system is not an author.

\end{document}
